\section{Conclusion and Future Outlook}
\label{sec:conclusion}

In this work, we challenged the fundamental and ubiquitous assumption that activation representations and task boundaries in modular deep learning are best modeled in flat Euclidean space ($\mathbb{R}^D$). We demonstrated that this Euclidean assumption introduces severe representation crowding and destructive inter-task cross-talk near the coordinate origin, leading to substantial accuracy degradation when multiple task-specific expert adapters are ensembled.

To shatter these geometric limits, we introduced \textbf{HyperMerge} (\textbf{Hyperbolic Space Activation Routing and Fusion}), a novel paradigm that shifts the geometric substrate of dynamic model ensembling from flat vector planes to the Poincaré Ball model of Hyperbolic Space ($\mathbb{D}_c^D$). By leveraging negative curvature, HyperMerge provides exponential volume expansion that naturally segregates disjoint task manifolds near the boundary, thereby completely resolving representation crowding. Through mathematically rigorous, closed-form formulations of \emph{Hyperbolic Centroid Alignment} (HCA) and \emph{Beltrami-Klein Symmetric Blending} (BKSB), HyperMerge performs single-pass, sample-wise non-linear ensembling with zero buffering latency and absolute immunity to stream heterogeneity. 

Empirical evaluation on a 14-layer Analytical Coordinate Sandbox demonstrated that HyperMerge achieves a flatline joint mean accuracy of \textbf{83.40\% $\pm$ 5.15\%} under both homogeneous and heterogeneous streams, completely eliminating representation collapse, operating safely within the logical expert ceiling of \textbf{99.98\% $\pm$ 0.02\%}, and performing on par with the state-of-the-art Euclidean ensembling scheme (SPS-ZCA) with a permutation-invariant formulation.

\textbf{Future Outlook:} 
While this work demonstrates the immense power of hyperbolic geometry in a simulated coordination sandbox, we believe we have only scratched the surface of non-Euclidean modular deep learning. In future work, we plan to scale HyperMerge to massive-scale foundation models, including Large Language Models (e.g., LLaMA, BLOOM) and Vision-Language Models (e.g., CLIP, LLaVA). Furthermore, exploring other non-Euclidean spaces, such as spherical manifolds for periodic task features, or product manifolds (combining hyperbolic, spherical, and Euclidean spaces) to adapt to mixed hierarchical and flat geometries, offers a rich and highly promising path forward. Ultimately, this work serves as an invitation to the community to look beyond the flatlands of standard Euclidean coordinates and explore the rich mathematical landscapes of Riemannian geometry to build more robust, expressive, and versatile modular systems.
